% Part:sets-functions-relations
% Chapter: size-of-sets
% Section: enumerations-alt

\documentclass[../../../include/open-logic-section]{subfiles}

\begin{document}

\olfileid{sfr}{siz}{enm-alt}

\olsection{प्रगणने आणि \usetoken{S}{enumerable} संच}

\begin{editorial}
  या विभागात संचसिद्धान्तात रूढ असलेल्या पद्धतीने, $\Nat$ च्या
  (आरंभीच्या खंडां)शी असलेली एकास-एक आच्छादने म्हणून प्रगणनांची
  व्याख्या केली आहे. त्यामुळे \olref[enm]{sec} मधील व्याख्यांशी तिचा
  किंचित विरोध होतो, आणि तेथील सर्व उदाहरणे येथे पुन्हा येतात. हा
  विभाग त्या विभागापेक्षा थोडा अधिक संक्षिप्तही आहे.
\end{editorial}

सांत संच त्यातील !!{element}s केवळ क्रमाने मांडून सांगता येतो. पुढीलप्रमाणे
संचाची व्याख्या करताना आपण हेच करतो:
\[
  A = \{a_1, a_2, \ldots, a_n\}.
\]
!!{element}s $a_1$, \dots, $a_n$ हे सर्व भिन्न आहेत असे गृहीत धरल्यास,
यातून $A$ आणि पहिल्या $n$ नैसर्गिक संख्या $0$, \dots, $n-1$ यांच्यात
!!a{bijection} मिळते. उलट, प्रत्येक सांत संचात सांत संख्येनेच
!!{element}s असल्यामुळे प्रत्येक सांत संच अशी संगती लावून मांडता येतो.
दुसऱ्या शब्दांत, $A$ सांत असेल, तर $A$ आणि $\{0, \dots, n-1\}$
यांच्यात !!a{bijection} असते; येथे~$A$ मध्ये $n$ इतके !!{element}s
असतात.

विशिष्ट प्रकारचे अनंत संच मान्य केल्यास, काही अनंत संचांचेही प्रगणन
करता येईल. अनंत संच आणि संपूर्ण~$\Nat$ यांच्यातील !!a{bijection} त्या
संचाचे प्रगणन करते, असे म्हणून हे नेमकेपणाने मांडता येते.

\begin{defn}[प्रगणनाची संचसैद्धान्तिक व्याख्या]
संच $A$ चे \emph{प्रगणन} म्हणजे असे !!a{bijection} की ज्याची व्याप्ती
$A$ असते आणि ज्याचा प्रांत नैसर्गिक संख्यांचा एखादा आरंभीचा खंड $\{0,
1, \ldots, n\}$ {किंवा} नैसर्गिक संख्यांचा संपूर्ण संच~$\Nat$ असतो.
\end{defn}

\begin{explain}
\emph{प्रगणन} या शब्दाच्या अशा वापरामागे सहज समजणारा आधार आहे.
संच $A$ चे प्रगणन केले आहे असे म्हणणे म्हणजे संच~$A$ मधील घटक मोजत
जाण्यास अनुमती देणारे !!a{bijection} $f$ आहे असे म्हणणे. $0$ वा घटक
$f(0)$, पहिला $f(1)$, \ldots आणि $n$ वा $f(n)$,~\ldots असतो.
\footnote{होय, आपण $0$ पासून मोजतो. अर्थात~$1$ पासूनही सुरुवात करता
येईल. त्यामुळे मोठा फरक पडणार नाही; फक्त~$\Nat$ ऐवजी~$\PosInt$
घ्यावा लागेल.} पुढील व्याख्या जोडल्यावर या मागचे कारण आणखी स्पष्ट होते:
\end{explain}

\begin{defn}
  \ollabel{defn:enumerable}
  संच~$A$ हा !!{enumerable} असतो तेव्हा आणि केवळ तेव्हाच, जेव्हा
  $A = \emptyset$ असते किंवा~$A$ चे प्रगणन असते. $A$ हा
  !!{nonenumerable} आहे असे म्हणतो तेव्हा आणि केवळ तेव्हाच, जेव्हा
  $A$ हा !!{enumerable} नसतो.
\end{defn}

\begin{explain}
म्हणून संच रिक्त असेल किंवा प्रगणन वापरून त्यातील !!{element}s मोजता
येत असतील तेव्हा आणि केवळ तेव्हाच तो !!{enumerable} असतो.
\end{explain}

\begin{ex}
नैसर्गिक संख्यांचे प्रगणन करणारे फलन म्हणजे फक्त $\Id{\Nat} \colon
\Nat \to \Nat$ हे $\Id{\Nat}(n) = n$ ने दिलेले एकरूपता फलन.
\emph{धन} नैसर्गिक संख्या $\Nat^+ = \Nat \setminus \{0\}$ यांचे
प्रगणन करणारे फलन $g(n) = n + 1$ आहे; म्हणजेच ते उत्तरवर्ती फलन आहे.
\end{ex}

\begin{prob}
संच $A$ हा !!{enumerable} असतो तेव्हा आणि केवळ तेव्हाच $A =
\emptyset$ असते किंवा $f\colon \Nat \to A$ असे !!a{surjection}
असते, हे दाखवा. $A$ हा !!{enumerable} असतो तेव्हा आणि केवळ तेव्हाच
$g\colon A \to \Nat$ असे !!a{injection} असते, हेही दाखवा.
\end{prob}

\begin{ex}
पुढीलप्रमाणे दिलेली $f\colon \Nat \to \Nat$ आणि $g \colon \Nat
\to \Nat$ ही फलने
\begin{align*}
  f(n) & = 2n \text{ आणि}\\
  g(n) & = 2n+1
\end{align*}
अनुक्रमे सम नैसर्गिक संख्यांचे आणि विषम नैसर्गिक संख्यांचे प्रगणन
करतात. पण यांपैकी कोणतेही फलन !!{surjective} नाही; म्हणून कोणतेही
$\Nat$ चे प्रगणन नाही.
\end{ex}

\begin{prob}
वर्गसंख्या $1$, $4$, $9$, $16$, \dots{} यांचे प्रगणन परिभाषित करा.
\end{prob}

\begin{ex}
$\lceil x \rceil$ हे $x$ ला त्याच्यापेक्षा कमी नसलेल्या सर्वांत जवळच्या
पूर्णांकाकडे नेणारे \emph{ऊर्ध्व पूर्णांक} फलन असू द्या. मग पुढीलप्रमाणे
दिलेले फलन $f \colon \Nat \to \Int$:
\[
  f(n) = (-1)^{n} \left\lceil\tfrac{n}{2}\right\rceil
\]
पूर्णांकांचा संच~$\Int$ पुढीलप्रमाणे प्रगणित करते:
\[
\begin{array}{c c c c c c c c}
f(0) & f(1) & f(2) & f(3) & f(4) & f(5) & f(6) & \dots \\ \\
\big\lceil \tfrac{0}{2} \big\rceil & -\big\lceil \tfrac{1}{2}\big\rceil &  \big\lceil \tfrac{2}{2} \big\rceil & -\big\lceil \tfrac{3}{2} \big\rceil & \big\lceil \tfrac{4}{2} \big\rceil  & -\big\lceil \tfrac{5}{2}\big\rceil & \big\lceil \tfrac{6}{2} \big\rceil & \dots \\ \\
0 & -1 & 1 & -2 & 2 & -3 & 3& \dots
\end{array}
\]
धन आणि ऋण पूर्णांकांमध्ये आलटूनपालटून ``उड्या मारून'' $f$ हे $\Int$
ची मूल्ये कशी निर्माण करते ते पाहा. $f$ ची व्याख्या पुढीलप्रमाणे
प्रकरणे पाडून केली आहे असेही मानता येते:
\[
f(n) = \begin{cases}
  \frac{n}{2} & \text{जर $n$ सम असेल}\\
  -\frac{n+1}{2} & \text{जर $n$ विषम असेल}
  \end{cases}
\]
\end{ex}

\begin{prob}
$A$ आणि $B$ हे !!{enumerable} असतील, तर $A \cup B$ सुद्धा गणनीय
आहे, हे दाखवा.
\end{prob}

\begin{prob}
$A_1$, $A_2$, \dots, $A_n$ हे सर्व !!{enumerable} असतील, तर
$A_1 \cup \dots \cup A_n$ सुद्धा गणनीय आहे, हे $n$ वरील गणिती विगमनाने
दाखवा.
\end{prob}

\end{document}
